%
% IEEE Transactions on Microwave Theory and Techniques example
% Tibault Reveyrand - http://www.microwave.fr
%
% http://www.microwave.fr/LaTeX.html
% ---------------------------------------



% ================================================
% Please HIGHLIGHT the new inputs such like this :
% Text :
%  \hl{comment}
% Aligned Eq. 
% \begin{shaded}
% \end{shaded}
% ================================================



\documentclass[journal]{IEEEtran}

% \documentclass[12pt, letterpaper,
% twoside]{article}

% \usepackage{emoji}
\usepackage[graphicx]{realboxes}
%\usepackage[retainorgcmds]{IEEEtrantools}
%\usepackage{bibentry}  
\usepackage{xcolor,soul,framed} %,caption
\usepackage{cite}
\usepackage{multirow}
\usepackage{fontawesome}
\colorlet{shadecolor}{yellow}
% \usepackage{color,soul}
% \usepackage[pdftex]{graphicx}
% \graphicspath{{../pdf/}{../jpeg/}}
\DeclareGraphicsExtensions{.pdf,.jpeg,.png}
\usepackage{caption}
\usepackage{subcaption}
\usepackage[cmex10]{amsmath}
\usepackage[colorinlistoftodos]{todonotes}
\usepackage{hyperref} % For clickable ORCID link
\usepackage{academicons} % For ORCID icon
\usepackage{xcolor} % For colors
\newcommand{\orcidicon}{
    \begin{tikzpicture}
    \draw[white, fill=white] (0,0) rectangle (1ex,1ex);
    \fill[green!60!black] (0, 0) rectangle (1ex, 1ex);
    \draw[white] (0.5ex, 0.5ex) circle (0.4ex);
    \draw[white, line width=0.07ex] (0.5ex, 0.5ex) -- ++(60:0.3ex);
    \draw[white, line width=0.07ex] (0.5ex, 0.5ex) -- ++(0:0.3ex);
    \draw[white, line width=0.07ex] (0.5ex, 0.5ex) -- ++(-60:0.3ex);
    \end{tikzpicture}
}
%Mathabx do not work on ScribTex => Removed
%\usepackage{mathabx}
% \usepackage{algorithm}
% \usepackage{algpseudocode}
\usepackage{comment}
% \let\hl[1]\includecomment{#1}
\usepackage{enumerate}% http://ctan.org/pkg/enumerate
% \newcommand{\hlcyan}[1]{{\sethlcolor{cyan}\hl{#1}}}
\newcommand\textblue[1]{\textcolor{blue}{#1}}
\newcommand\textred[1]{\textcolor{red}{#1}}
\definecolor{orcidgreen}{HTML}{A6CE39}

% \newcommand\myworries[1]{\textcolor{red}{#1}}
\newcommand{\orcid}[1]{\href{https://orcid.org/#1}{\textcolor{orcidgreen}{\aiOrcid}}}

\newcommand\quickthings[1]{\textblue{\\\faQuestion #1}}
\newcommand\maybelater[1]{\textred{\\\faClockO#1}}
% \newcommand\quickthings[1]{\textblue{\\\faQuestion #1}}
\newcommand\peter[1]{\textcolor{red}{\faComment #1}}
\newcommand\abs[1]{\\\hl{#1}}



%Comment to keep todo parts, Uncomment to remove TODO parts!
\renewcommand\quickthings[1]{}
% \renewcommand\absolutelynecessary[1]{}
\renewcommand\peter[1]{}
\renewcommand\maybelater[1]{}
\renewcommand\abs[1]{}
\renewcommand\hl[1]{#1}

\usepackage[linesnumbered,ruled,vlined]{algorithm2e}



% % \usepackage
% \usepackage[ruled,vlined]{algorithm2e}
% \DeclareMathOperator{\ind}{ind}
% \DeclareMathOperator{\val}{val}
% \DeclareMathOperator{\ptr}{ptr}
% \DeclareMathOperator{\row}{row}
% \usepackage{algpseudocode}
\hyphenation{op-tical net-works semi-conduc-tor}

%\bstctlcite{IEEE:BSTcontrol}


%=== TITLE & AUTHORS ====================================================================
\begin{document}
\bstctlcite{IEEEexample:BSTcontrol}
    % \title{Adversarial attacks on federarted learning for medical image analysis}
    \title{ A Comparative Study of Federated Learning Methods for COVID-19 Detection  }
 % \author{Michael~Shell,~\IEEEmembership{Member,~IEEE,}
 %        John~Doe,~\IEEEmembership{Fellow,~OSA,}
 %        and~Jane~Doe,~\IEEEmembership{Life~Fellow,~IEEE}
  
  % \author{Erfan~Darzi\href{https://orcid.org/0000-0002-6003-0119}{\includegraphics[height=1.7ex]{orcid}} ,~\IEEEmembership{Member,~IEEE},
  % Nanna~M.~Sijtsema\href{https://orcid.org/0000-0001-6644-274X}{\includegraphics[height=1.7ex]{orcid}}, and~P.M.A~van~Ooijen\href{https://orcid.org/0000-0002-8995-1210}{\includegraphics[height=1.7ex]{orcid}},~\IEEEmembership{Senior Member,~IEEE}
 

    %   and~Zoya~Popovi\'c,~\IEEEmembership{Fellow,~IEEE}% <-this % stops a space
% \\\textit{$^{1}$Machine learning lab, Data Science Center in Health (DASH) \\University of Groningen, Hanzeplein 1, Groningen, The Netherlands}
% \\ e.darzidehkalani@rug.nl
%   \thanks{ The work of Erfan Darzi was supported by KWF Kankerbestrijding and the Netherlands Organisation for Scientific Research (NWO)  Domain AES, as part of their joint strategic research programme: Technology for Oncology IL. The collaboration project is co-funded by the PPP allowance made available by Health Holland, Top Sector Life Sciences \& Health, to stimulate public-private partnerships. (\textit{Corresponding author: Erfan Darizdehkalani})\\
%   Erfan Darizdehkalani, Nanna M. Sijtsema, and P.M.A van Ooijen are with the Machine learning lab, Data Science Center in Health (DASH)
% University of Groningen, Hanzeplein 1, Groningen, The Netherlands.  (e-mail: e.darzidehkalani@rug.nl, n.m.sijtsema@umcg.nl, p.m.a.van.ooijen@umcg.nl)
  
% }
  % \thanks{Erfan Darzidehkalani is with UMCG (e-mail: e.darzidehkalani@umcg.nl).}% <-this % stops a space

% The paper headers
% \markboth{IEEE TRANSACTIONS ON MICROWAVE THEORY AND TECHNIQUES, VOL.~60, NO.~12, DECEMBER~2012
% }{Roberg \MakeLowercase{\textit{et al.}}: Cross modality image transformation using cyclic residual
% architecture and evolutionary algorithm}
% ===================================================================
\maketitle
% === ABSTRACT ====================================================================
% =================================================================================
\begin{abstract}

Deep learning has proven to be highly effective in diagnosing COVID-19; however, its efficacy is contingent upon the availability of extensive data for model training. The data sharing among hospitals, which is crucial for training robust models, is often restricted by privacy regulations. Federated learning (FL) emerges as a solution by enabling model training across multiple hospitals while preserving data privacy. However, the deployment of FL can be resource-intensive, necessitating efficient utilization of computational and network resources. In this study, we evaluate the performance and resource efficiency of five FL algorithms in the context of COVID-19 detection using Convolutional Neural Networks (CNNs) in a decentralized setting. The evaluation involves varying the number of participating entities, the number of federated rounds, and the selection algorithms. Our findings indicate that the Cyclic Weight Transfer algorithm exhibits superior performance, particularly when the number of participating hospitals is limited. These insights hold practical implications for the deployment of FL algorithms in COVID-19 detection and broader medical image analysis.
\end{abstract}
% \section{Background and Related works}

 % $\Delta W_n^t$
 
Federated learning has demonstrated efficacy in an array of imaging modalities, including Magnetic Resonance Imaging (MRI) \cite{sheller2020federated}\cite{silva2019federated}, X-ray \cite{balachandar2020accounting}, retinal imaging \cite{balachandar2020accounting}, as well as in applications such as brain tumor segmentation \cite{bakas2017advancing}\cite{lee2018privacy}, diagnosis \cite{pan2019improving}, and treatment selection \cite{lee2018privacy}. In particular, Federated Learning (FL) has proven to be a valuable tool for supporting physicians in their decision-making process regarding the treatment of COVID-19 patients. A landmark study that involved 20 institutions across five continents found that FL played a significant role in shaping patient treatment plans\cite{flores2021federated}. The study employed chest radiography images in conjunction with clinical information to determine the appropriate level of care and oxygen requirements for patients afflicted with COVID-19. It was observed that FL improved the performance of the predictive model, especially for institutions with smaller datasets, compared to using only local data for model training. Additionally, it was found that healthcare facilities with smaller datasets often had underrepresented categories due to a low number of patients in certain classes. The implementation of FL led to a notable improvement in predictions for these underrepresented patient categories.

\begin{figure*}[t!]
\centering
\includegraphics[width=1\textwidth]{FLmodels.png}
\caption{Illustration of FL models and algorithms: (a) Federated averaging, where clients train on a local batch of data. (b) FedSGD, in which a subset of clients is selected, and each performs a single step of SGD before sending model updates to the server. (c) Cyclic Weight Transfer (CWT), where clients train locally and pass the model to the next client, repeating the cycle. (d) Single Weight Transfer (SWT), where the model passes through each client only once. (e) Stochastic Weight Transfer (STWT), in which the model is sequentially passed through clients, with participating clients in each round being sampled randomly.}
\label{fig:flalgorithms}
\end{figure*}


Recent research has focused on the classification of scan images to distinguish between COVID-19 patients and healthy individuals, as well as identifying lesion areas. The primary application of AI in managing COVID-19 patients has been the interpretation of radiology images, especially chest CT scans.  The detection of lung alterations through these scans plays an important role in optimizing patient management and guiding treatment decisions\cite{yan2020interpretable}\cite{hu2020challenges}\cite{burian2020intensive}.  Several studies have also explored 3D Convolutional neural networks \cite{wang2020weakly} and COVID-19 detection with a limited number of training samples.

While the majority of these studies report favorable accuracy, they often presume a centralized environment wherein a single data center has access to all data. However, a few studies have successfully applied distributed learning for COVID-19 detection, employing global aggregation models such as model averaging in federated learning settings\cite{ho2022fedsgdcovid}\cite{zhang2021dynamic}, or within a blockchain framework \cite{kumar2021blockchain}. These studies have pointed to certain limitations of existing algorithms, such as high communication overhead \cite{remedios2020federated}, as well as convergence issues or catastrophic forgetting when the number of participating hospitals increases \cite{sheller2020federated} \cite{chang2018distributed}.


To our knowledge, no study has been performed that compared multiple FL algorithms under standard conditions to evaluate their applicability. Therefore, comparing multiple FL algorithms under standard conditions could be informative in evaluating their applicability in practice.

To evaluate the existing methods from multiple perspectives, we have implemented the most popular models and compared them in terms of performance, communication overhead, and computation burden.

% 
\section{Algorithms}
% \label{sec:algorithms}
\textbf{Centralized Data Sharing} In Centralized Data Sharing (CDS), data is stored in a central location and is accessible to all clients. This stands in contrast to federated and decentralized data sharing methods, where data is stored across multiple locations and accessed by either a single user or a limited number of users. CDS serves as a baseline for comparing other algorithms.
\textbf{Federated Averaging} Federated Averaging involves an iterative learning procedure comprising local and global steps. In this process, each data owner trains a model received from a global server on its local dataset through local iterations \cite{mcmahan2017communication}. Subsequently, the global server aggregates the updated local models to update the global model. This global model is then distributed to clients for the subsequent round. The optimization problem for Federated Averaging can be expressed as:
\begin{equation}
w^{t+1} = \sum\limits_{i=1}^{N}{p_{i} w_{i}^{t}} , w_{i}^{t}=\arg\min\limits_{w_{i}}{\left(\mathcal{L}(\mathcal{D}_{i};w^{t})\right)}
\end{equation}
where $N$ is the number of data owners, $\mathcal{L}(\mathcal{D}_{i};w^{t})$ is a loss function indicating global model parameters  $w^{t}$ of local datasets, and $p_{i}$ is the probability of selecting client $i$. 
Local optimization can be formulated as $w_{i}^{t+1} \leftarrow w^{t}-\eta\cdot \nabla \mathcal{L}(w^{t};\mathcal{D}_{i})$, where 
 $\eta$ is the learning rate. The global model can be updated based on the local models $w_{i}$ and is shared for aggregation: 
\begin{equation}
w^{t+1} = \sum\limits_{i=1}^{N}{p_{i} w_{i}^{t+1}}
\end{equation}

\textbf{Federated Stochastic Gradient Descent} Federated Stochastic Gradient Descent, or FedSGD\cite{chai2020fedeval}, is a variant of Federated Averaging (FedAvg) that employs a large-batch synchronous approach for multi-client learning. In FedSGD, a subset of clients is chosen from the total pool, with this subset being defined by $C$. This selection is made at each global round. Subsequently, the global server dispatches the latest global model to the chosen clients. Each of these clients then undertakes local training on its dataset for a predetermined number of epochs. The global model is updated in light of the local models obtained from each client and is made available for aggregation, a process that is analogous to that in FedAvg. Nevertheless, FedSGD distinguishes itself by computing the gradient over the chosen batch of clients, hence $C < 1$. When $C = 1$, the training becomes non-stochastic, or full batch, as it involves all the clients. This makes it possible to train with large batches since the gradient is computed over the selected subset of clients.

The optimization problem for FedSGD can be expressed as follows:
\begin{equation}
w^{t+1} = w^{t} - \eta \cdot \sum\limits_{i=1}^{C}p_i \cdot \nabla \mathcal{L} (w^t;\mathcal{D}_i)
\end{equation}
where $\eta$ is the learning rate, $p_i$ is the probability of selecting client $i$ and $\mathcal{L}$ is the loss function.
The key difference between FedAvg and FedSGD lies in the use of large-batch synchronous approach in FedSGD. This approach has been shown to outperform the naive asynchronous SGD training due to the increased accuracy and efficiency, as compared to the local training approach used in FedAvg \cite{chai2020fedeval}\cite{charles2021large}. Additionally, FedSGD has been shown to be more robust to non-IID data distributions, compared to FedAvg \cite{chai2020fedeval}. 
\quickthings{This part is unclear to me.}
% \absolutelynecessary{}
\maybelater{paraphrised}

\textbf{Cyclic weight transfer} Federated learning techniques have been widely used in medical image processing tasks using a method known as cyclic weight transfer (CWT)\cite{balachandar2020accounting}. This method involves training models on individual clients for a number of iterations and then cyclically sharing the updated weights with the following client. However, the existing CWT algorithm faces a notable challenge, as it lacks the ability to effectively manage inter-client variability in training data or labels.  To ensure the practical application of CWT, it is crucial to develop a version that can handle the common variations observed in a majority of real-world medical imaging datasets\cite{darzidehkalanifederatedII}.

\textbf{Single weight transfer} Single weight transfer (SWT) is anothe FL method widely used in the medical imaging domain. In Single weight transfer, models are trained in each client with its local data, and then the updated model is transferred to the next client. The difference between this method and CWT is that here the model passes each client only once\cite{darzidehkalanifederatedI}. 

\textbf{Stochastic weight transfer}
In stochastic weight transfer (STWT), we select a subsample of clients and train them in a cyclic manner. Similar to FedSGD, a ratio defines the number of selected clients to the total number of clients in each federated round. 

Figure \ref{fig:flalgorithms} provides a visual representation of the introduced algorithms.



% سوال / مسیله چیه
% Importance: Why your research matters in the context of an industry or the world
% اهمیت موضوع
% 	کجاها کاربرد داره
% مشکلات موضوعات قبلی و اینکه مسایل دگ کار نمیکنن
% برای اینکه این مسایل کار کنن نیازی به چه چیزی بود. ک اونا کم داشتن
% چرا این روش تو داره کار میکنه
% چه سوالی رو جواب میده
% توضیح اینکه چجوری کار میکنه
% تعریف و تمجید از کار کردنش

% \maybelater{chanta akse tamiz ba coreldraw bekesh}

% Their
% We are comparing various implementations of FL and comparing their performance level. The performance is done by metrics.
% 2. Although FL has been introduced to tackle the problem of privacy, its distributed nature requires too much commuincation between clients. Knowsing can determine its deployability and scailibitlty. Hence the question is how to these model compare it terms of communication between clients.
% 3. Computation time. Each FL model is investigated based on the required time for computation.
% 4. Effect of rounds: Number of Federated rounds 


% \quickthings{Maybe you can add a benchmarking section similar to "A Performance Evaluation of FL Algorithms"}
% \maybelater{Some ideas are herer https://arxiv.org/pdf/1709.05929.pdf
% }



% 

\section{Experiment}




\subsection{Datasets}

\textbf{Brain cancer classification:}Dataset for detecting brain cancer is downloaded from Kaggle. It contains brain MRI images categorized into four classes: Three types of brain tumor,  and one healthy no Tumor class. 
% For tumor type classification task, 2870 images for training and 394 images for the test. 
For Meningioma and Glioma detection, 1437 and 1426 samples were used, respectively,  and they were split as train/ test data. Transforms rotation, Flipping, and normalizing.  Images were resized to 100$\times$100. 
\\
\textbf{Histopathologic cancer detection:} 
The samples are metastasic tissue images of lymph node cancer. Images are in size ,96 $\times$ 96, and the task is to classify the tissue samples. Each positive sample has a metastasic region, which is located in $32\times32$ neighborhood of each image sample. 
% Tumor tissue in the outer region of the patch does not influence the label. 
They were categorized into cancer and non-cancer classes. 2150 samples were chosen and randomly assigned to clients, and each split 62\% for train and 38\% for the test. We used horizontal Flipping, and images were normalized.

% \\\hl{A ro tu line 21 algorithm taarif kon va dorost kon notionesh ro}
\subsection{Network architectures}

\textbf{Deep learning models:}
The deep learning model used is a Convolutional Neural Network (CNN) with six layers of convolution stacked to 5 fully connected layers. The activation function used is ReLU. And dropout parameter (0.25). The last layer depends on the number of classes (4 or 2). Models are all trained and converged before implementing adversarial attacks.
They are trained with Cross-Entropy loss and with an SGD optimizer. Four models are trained separately. Two binary classification models are trained to detect Meningioma and Glioma from the healthy No Tumor class. One model is trained to detect cancer in HistoPathology images.

\textbf{Federated setting:}
Three clients are defined for the FL setup. The data is assigned randomly, and the clients have non-IID data distribution.
The FedAVG method is used to aggregate the models. The aggregation is weighed based on the length of the training dataset. Each client is trained for 20 epochs at the communication round. The total FL rounds are 50.

\subsection{Attack setting}

Each client is assigned 100 test samples to perform the adversarial attack. Each client is selected as the adversary, and the results are averaged at the end. 
% The attacks are conducted with different $\epsilon$ and gradient steps $\alpha$ to evaluate the effect on the final result. 
We use three metrics,
Clean accuracy is defined as the performance of models on uncorrupted test images, \textit{Attack Success Rate (ASR)}; how much an adversary can change the predicted labels produced by each model, \textit{Average Error Transferability Rate (AETR)} to evaluate transferability. An adversary might propose an attack by only selecting the samples that successfully fooled its own model. AETR measures how successfully these samples changed the target models' correct predictions. So unlike ASR, this metric doesn't count the examples that were initially misclassified by the victim model.\cite{papernot2017practical},

\begin{figure*}[h!]
     \centering
     \begin{subfigure}
         \centering
         \includegraphics[width=0.32\textwidth]{Meningioma_ASR_EPS.eps}
          \includegraphics[width=0.32\textwidth]{Glioma_ASR_EPS.eps}
           \includegraphics[width=0.32\textwidth]{Pathology_ASR_EPS.eps}
         \caption{Effect of $\epsilon$ on Average ASR on benign clients }
         \label{fig:AASR benign-EPS}
     \end{subfigure}
     \hfill
     \begin{subfigure}
          \centering
         \includegraphics[width=0.32\textwidth]{Meningioma_others_EPS.eps}
          \includegraphics[width=0.32\textwidth]{Glioma_others_EPS.eps}
           \includegraphics[width=0.32\textwidth]{Pathology_others_EPS.eps}
         \caption{Effect of $\epsilon$ on ASR on the advesarial client }
         \label{fig:AASR malignant-EPS}
     \end{subfigure}
     \hfill
        \caption{Effect of Error perturbation degree $\epsilon$ on attack transferability. FGSM and PGD attack with and without CRN initialization where performed. CRN initalizations used 10 models and another one loading 5.  ASR is calculated on benign and adversary clients. The results of with varying perturbation degrees. The higher ASR on benign clients shows higher transferability}
        \label{fig: epsilon graphs}
\end{figure*}

  Here, we expand the previous findings to the FL settings and discuss whether change in $\epsilon$ can lead to transferability. 
  
We perform three analysis to find important factors in attack success:
\begin{itemize}
    \item We investigate dependency of attacker on $\epsilon$. By visual inspection and ASR evaluation, and as a comparison to previously found optimal known values in centralized  setting.
    \item  We compare  efficiency of models, by discussing their attack preparation time and their ASR.
    \item We also see how attack step $\alpha$ can determine ASR.
\end{itemize}



% \section{Results}







% Original images, adversarial images, and corresponding adversarial noise created with FGSM ( ϵ=0.02) in different black-box settings


\subsection{Efficiency analysis}


\begin{figure*}[h!]
     \centering
     \begin{subfigure}
         \centering
         \includegraphics[width=0.32\textwidth]{Meningioma_ASR_EPS_steps.eps}
          \includegraphics[width=0.32\textwidth]{Glioma_ASR_EPS_steps.eps}
           \includegraphics[width=0.32\textwidth]{Pathology_ASR_EPS_steps.eps}
         \caption{Average ASR on benign clients }
         \label{fig:y equals x}
     \end{subfigure}
     \hfill
     \begin{subfigure}
          \centering
         \includegraphics[width=0.32\textwidth]{Meningioma_others_EPS_steps.eps}
          \includegraphics[width=0.32\textwidth]{Glioma_others_EPS_steps.eps}
           \includegraphics[width=0.32\textwidth]{Pathology_others_EPS_steps.eps}
         \caption{Average ASR on advesarial client }
         \label{fig:y equals x}
     \end{subfigure}
     \hfill
        \caption{Effect of Error perturbation step $\alpha$ on attack transferability. PGD attack with CRN initalizations used 10 models and another one loading 5. ASR is calculated on benign and adversary clients. The results of with varying perturbation steps. The higher ASR on benign clients shows higher transferability}
        \label{fig:three graphs}
\end{figure*}

The models are trained in an FL enviroment. Average test results, or clean accuracy shows the average performance of clients on unperturbed test data. This could be a baseline for comparison of attacks.
Here we compare the computational complexity of iterative algorithms. The comparison is based on the time required to train one batch of adversarial examples, with  $\epsilon=0.05$ and $\alpha=0.001$. The average Attack success rate is the average of ASR for all test samples of all clients and indicates the transferability in the FL network.
As expected \ref{table_time}, the higher iteration rounds take more computation and lead to higher AASR. Generally, computational complexity has a linear dependence on the number of iterations. However, the increase in transferability can vary. If PGD performs poorly initially, having higher iterations doesn't help much. 

Also, CRN does not require additional computation. Similar to single-step methods, computational load is small but increases on high-resolution images. As an example, the pathology dataset requires more calculation than MRI data. 
However, the overall time is far less than iterative models, and also it could enhance the transferability much better. The effect is consistent among different datasets. 
% However,  This means that Pathology images have better AASR than MRI images in baseline and CRN enabled attacks.
% Also,  

% PGD can sometimes perform poorly, even in higher iterations. Enabling the initialization might  


\begin{table}[h!]
\centering
\setlength{\tabcolsep}{6pt}
\renewcommand\arraystretch{1.22}
\caption{ \small Comparison of iterative models and their copmutational efficiency, on performing computations on one batch of data.  ACC shows average client performance for unperturbed test data. AASR is average ASR on all clients.}
\begin{tabular}{| *{5}{c|} }
\hline
Dataset  & Attack type & ACC & AASR & time (sec)
\\   \hline  
\multirow{3}{5em}{Meningioma}     &\textbf{PGD-1+CRN10}&\multirow{3}{3em}{84.12\%}&\textbf{63.92\%} & \textbf{0.217}  \\
 &PGD-20&&27.52\% & 3.423 \\
&PGD-40&&32.99\% & 6.794  \\ \hline
\multirow{3}{4em}{Pathology}     &\textbf{PGD-1+CRN10}&\multirow{3}{3em}{77.01\%}&\textbf{90.92\%} & \textbf{0.354} \\
&PGD-20&&60.98\% & 5.464\\
&PGD-40&&82.27\% & 10.843 \\ \hline
\multirow{3}{3em}{Glioma}     &\textbf{PGD-1+CRN10}&\multirow{3}{3em}{61.84\%}&\textbf{70.55}\% & \textbf{0.216}  \\
&PGD-20&&51.83\% &3.420  \\
&PGD-40&&63.77\% & 6.793  \\ \hline


\end{tabular}
\label{table_time} 
\end{table}




\subsection{Effect perturbation degree}

We performed attack scenarios to evaluate the effect of perturbation degree $\epsilon$  on attack success on the model. We compared the baseline attack to CRN- enabled attacks. ASR is calculated on benign and adversary clients and is shown in Fig \ref{fig: epsilon graphs}.

More successful attacks on adversaries led to higher transferability to benign clients. Low values of $\epsilon$ can decrease transferability to a large extent. $\epsilon$ = 0.01 has poor performance on all the clients. Generally, Higher $\epsilon$ values lead to higher ASR in all scenarios. Although for Meningioma and Glioma images, increasing $\epsilon$ from 0.04 didn't have much positive effect in PGD and FGSM methods on the adversarial client. 
In pathology images, Iterative models can reach perfect accuracy on the adversarial client and have very high transferability.

CRN  has a positive effect on transferability in all scenarios. Also, CRN-enabled models are more dependent on $\epsilon$; their extent can vary depending on the dataset, attack method, and value of $\epsilon$.
Tuning $\epsilon$  might be tricky since performance and imperceptibility should be considered together.High $\epsilon$ values mean more distortion in the image. We performed an FGSM attack with varying $\epsilon$ values to evaluate the perturbation effect visually. Fig \ref{fig:perturbedMRI} shows the results of perturbed images with different perturbation degrees and their comparison with the unperturbed image. Note that perturbation values below  $\epsilon=0.05$ cause limited distortion.


Iterative models can reach high ASR very quickly with increasing $\epsilon$. 

\begin{figure*}[h!]
 \centering
 \includegraphics[width=0.8\textwidth]{Adversarial_attacks-5.eps}
 \caption{\small From (a) to (f), normal tumor image, and petrurbed images with FGSM attack, with perturbation parameters: $\epsilon$ = 0.01, 0.03, 0.05,  0.10, respectively}
 \label{fig:perturbedMRI}
\end{figure*}











% \subsection{AATR}
% \begin{table}[h]
% \centering
% \setlength{\tabcolsep}{8pt}
% \renewcommand\arraystretch{1.4}
% \caption{Result of our experiments for PGD attack on Pathology data $\epsilon$ values. The iterative models are trained for 40 rounds with step size of 0.01. \hl{\faQuestion maybe you can change the word for without CRN}}
% % \begin{tabular}{cccccccccccc}
% \begin{tabular}{| *{4}{c|} }\hline
% % Dataset  & Attack type & \multicolumn{2}{c|}{ $\epsilon=0.01$} & \multicolumn{2}{c|}{ $\epsilon=0.03$}           &    \multicolumn{2}{c|}{ $\epsilon=0.05$}    & \multicolumn{2}{c|}{ $\epsilon=0.07$}                                                              \\ \hline
%  Dataset&\multicolumn{3}{c|}{ FGSM}    
% %  &              \multicolumn{3}{c|}{ PGD}                                                                           \\  

%  \cline{2-5}
% &Baseline&CRN5 &CRN10\\\hline

% Meningioma&84.80\%&\textbf{86.31}\%&{84.02\%}
% % \%&83.73\%  &\textbf{84.25}\%  &84.16\%
% \\\hline
% Glioma&82.32\% &87.06\%&\textbf{89.98\%}
% % &74.89\%&\textbf{83.64}\% &82.32\%
% \\\hline

% Pathology&\textbf{90.67\%}&{84.81\%}\%&{85.98\%}
% % \%&79.86\% &86.79\%&\textbf{86.93\%}
% \\\hline



% \end{tabular}
% \label{table_core} 
% \end{table}










\subsection{Effect perturbation step}


Another critical parameter is the step of change in iterative methods. {$\alpha$} is shown to be highly deterministic in specific settings, but is not yet investigated in MIA domain. To evaluate how $\alpha$ can affect transferability, we examined the clients with $\eps=0.03$. Under various steps. The general observation is that there is an optimal middle range the same as $\epsilon$. The step size of 0.007 led to the highest transferability in MRI images. Larger steps led to a sharp decrease in transferability. For pathology images, midspans, about 0.007 led to the best ASR on the adversary, and higher values decreased both ASR on adversary and transferability. However, smaller steps also had high ASR values. 







\subsection{Error tranferability}




Another measure is the average error transfer rate (AETR). Although PGD outperforms FGSM in ASR, in AETR, the results are close. PGD is trained for 40 iterations, and the models are compared with and without CRN.
Both FGSM and PGD having high AETR suggest that although FGSM might not be able to produce high-quality samples in general, the good examples which can fool the adversary are also highly transferable to benign clients. And can be better than PGD samples. CRN generally increases the AETR. 




\begin{table*}[h!]
\centering
\setlength{\tabcolsep}{8pt}
\renewcommand\arraystretch{1.4}
\caption{Result of Average Error Transfer Rate (AETR) for FGSM and PGD methods, with and without CRN initalization }
% \begin{tabular}{cccccccccccc}
\begin{tabular}{| *{7}{c|} }\hline
% Dataset  & Attack type & \multicolumn{2}{c|}{ $\epsilon=0.01$} & \multicolumn{2}{c|}{ $\epsilon=0.03$}           &    \multicolumn{2}{c|}{ $\epsilon=0.05$}    & \multicolumn{2}{c|}{ $\epsilon=0.07$}                                                              \\ \hline
 Dataset&\multicolumn{3}{c|}{ FGSM}     &              \multicolumn{3}{c|}{ PGD}                                                                           \\  

 \cline{2-7}
&Baseline&CRN5 &CRN10&Baeline&CRN5&CRN10\\\hline

Meningioma&84.80\%&\textbf{86.31}\%&{84.02}\%&83.73\%  &\textbf{84.25}\%  &84.16\%\\\hline
Glioma&82.32\% &87.06&\textbf{89.98}&74.89\%&\textbf{83.64}\% &82.32\%\\\hline

Pathology&\textbf{90.67\%}&{84.81}\%&{85.98}\%&79.86\% &86.79\%&\textbf{86.93\%}\\\hline



\end{tabular}
\label{table_core} 
\end{table*}


% % \newpage

% 
\section{Discussion}
\label{sec:discussion}

% Local training had near random accuracy, but in all the experiments, the performance results were improved. 
% Considering low number of samples for individual clients, a training process limited to few clients results in near random classification accuracy. 
% Local models can have low performance due to limited datasets.
% Also,, ImageNet pre-training had a speed-up effect on our model performance.

% \textbf{Fundamental groups of FL}

% \textbf{how many groups}
% \textbf{how can we group them, name different ways}
% \textbf{effect number of data}


% Paragraph 1: This paragraph provides a “big picture” perspective for readers to remind them of the importance of your study.

% \textbf{overall results}
% 
Our results show that FL has comparable performance to centralized data sharing, with the advantage of keeping data private. With large volumes of data and after high number of rounds, centralized data sharing and cyclic weight transfer have the highest accuracy. 


Sequential models are susceptible to catastrophic forgetting, where a global model performs well on the latest client it has seen while having poor performance in other clients.\quickthings{cryptic sentence}\maybelater{I have paraphrised the sentence}
Conversely, in algorithms like FedAvg and FedSGD, the models are averaged asynchronously after all the clients have finished their training. So the trajectory is smoother and overall improving with more communication rounds. As shown in Figure \ref{seq-vs-noseq}, local test results can have a high variance when passing through clients sequentially, indicating the catastrophic forgetting effect.

Models like FedAvg, and FedSGD, in which all the clients have identical copies of one global model, are slower and more challenging to converge compared to sequential models like CWT and STWT. Also, FedAvG and FedSGD require more training resources due to active server participation, resulting in more computation and network consumption. Stochastic client selection is an efficient way of training. Stochastic models save significant time and resources while having similar performance to full client participation. Overall, CWT and STWT have best results in terms of model accuracy and computation times. These findings could be practical in further federated deployments in medical institutions.
% but when there are communication and computation limitations, other methods might be more practical.




% As can be seen in the results, CWT and STWT have better local results than other methods.

% So each local model is different than other models, and even in fewer rounds they retrain the model on their local data and achieve high accuracies. 

%  Unsurprisingly, CDS has higher accuracy than non-sequential methods.
 
% so in limited number of rounds, latter algorithms are preferred.



%  FedAVG needs to perform both local training and global aggregation, which requires is to have more computation, than sequential models which eliminate global aggregation and broadcasting step.   
%  Models with global server, suh as FedAVG and FedSGD, required more data transfer, because of the two way communication of each client with server.
Sequential models like CWT and STWT perform better than non-sequential models on fewer training rounds. For example,  after three rounds of training, STWT and CWT both reach 96\% accuracy, while FedAvG reaches 66\%, and FedSGD performs equally to a random classifier. As the training proceeds, FedAVG and FedSGD gradually improve with more global rounds.The concept of sequential models is similar to fine-tuning \cite{chen2020online} in centralized deep learning, so in cases where a hospital temporarily joins an FL network, or there is an urgency in training, sequential models are a better option.

More training rounds do not always lead to a better global model. Although average performance on all clients improves, more global rounds lead to worse performance for some clients. The global model can overfit some clients, leading to lower performance on others\cite{mohri2019agnostic}. Some studies suggested early stopping and fine-tuning to local dataset after global training is finished \cite{yu2020salvaging}. 
In all the algorithms, more clients resulted in slower convergence. This effect is stronger in the FedAvg algorithm. In FedAVG, the Global model must compromise between potentially disparate local minima\cite{li2019convergence}. Methods such as adaptive or stochastic selection of clients and momentum-based models help faster convergence \cite{liu2020accelerating}.
Our results suggest that stochastic client participation is close to full client participation. The average results of four trials with varying rounds, shown in Table \ref{number_of_rounds} indicate that stochastic client participation in FedSGD results in 5.23\% performance loss and 40\% less bandwidth consumption compared to FedAvg. In STWT, it results in only 1.25\% less accuracy but saves 40\% of communication and 11.3\% of computation.
% Also, SWT improves local clients' performance up to 80\% accuracy with 8 clients, and it requires extremely low bandwidth requirements.
These results are in accordance with prior studies, showing that, in theory, stochastic and full client participation have similar global minima\cite{cho2020client}. Stochastic client selection can be advantageous when there are limited resources, or in larger networks with occasionally unavailable clients.


% Models like FedSgd and STWT, select the participating clients in a stochastic manner. 
% The reason is because each client has less data, and there will be more number distant local client minima leading to unstable global model.


% Such can be a reflection of real world setting, where not all the clients are active in each round. So a fraction are selected.


% Some models have high max test accuracy while average is low indicating bias and catastrophic forgetting. Infrastructure limitations should be considered, when a central server has limited computing power it would be unable to perform FedAVG, instead no-server algorithms require to pass the model to the next client.

We did not assume any shift in clients' data. A more comprehensive analysis should consider the effect of the domain and distribution shifts on the performance of the algorithms. Also, inter-client data variability and the effect of heterogenous clients could be a future line of research.
% For deployments in longer timeframes, there might be changes in clients data distirbution , and previous results might be inapplicable Effect of domain shifts, and inter-client data variability could be a future line of research. 






\section{Conclusion}
\label{sec:conclusion}
 
FL enables extensive collaborations of hospitals to address medical imaging problems while keeping data private. Real-world implementation requires consideration of efficiency and hardware requirements in addition to model performance, especially in the healthcare field, which generally has limited infrastructure. We implemented five FL algorithms for COVID-19 detection and analyzed their efficiency and accuracy.
Our results suggest that FL algorithms have comparable performance to centralized data sharing, with the advantage of keeping data private. They also show that the sequential methods are a better option in most of the scenarios. This study can be helpful in the deployment of FL systems in COVID-19 detection and medical image analysis in general.

% demonstrate better performance for detecting COVID-19 patients and might be practical in deploying FL algorithms for covid-19 detection and medical image analysis in general.





\printbibliography







% === KEYWORDS
% ====================================================================
% =================================================================================
\begin{IEEEkeywords}
federated learning, medical image analysis, COVID-19, privacy preserving machine learning
\end{IEEEkeywords}






% For peer review papers, you can put extra information on the cover
% page as needed:
% \ifCLASSOPTIONpeerreview
% \begin{center} \bfseries EDICS Category: 3-BBND \end{center}
% \fi
%
% For peerreview papers, this IEEEtran command inserts a page break and
% creates the second title. It will be ignored for other modes.
\IEEEpeerreviewmaketitle


% ====================================================================
% ====================================================================
% ====================================================================










% \tableofcontents
% === I. INTRODUCTION =============================================================
% =================================================================================
\section{Introduction}

\IEEEPARstart{C}oronaviruses comprise a family of viruses known to cause respiratory and intestinal illnesses in both humans and animals. Among these viruses, the variants responsible for COVID-19, SARS, and MERS epidemics have gained notable attention. In cases of COVID-19, certain individuals might develop severe complications, including pneumonia, which can be detected through lung CT scans. Studies have indicated that chest imaging plays a crucial role in the diagnosis of COVID-19 in individuals exhibiting severe symptoms. In this domain, deep learning methods, especially Convolutional Neural Networks (CNNs), have proven to be highly effective in assisting radiologists in performing various image analysis tasks related to the diagnosis of COVID-19\cite{kogilavani2022covid}.

Deep learning models developed specifically for detecting COVID-19 infections have exhibited remarkable potential in identifying infected regions in CT scans and X-ray images. However, the training of these deep learning models necessitates access to ample and diverse medical datasets, which are typically collected from multiple sources. A majority of the current approaches rely on a centralized server for aggregating data from various healthcare institutions. This approach poses a challenge, as medical images often contain sensitive and confidential patient information, which is not permissible to share beyond the confines of the originating institution.

Federated learning (FL) emerges as a viable solution to this challenge by decentralizing the training process and retaining the data at its source. In a federated learning framework, distinct clients participate in the training process in a distributed manner using their local data. Specifically, each client independently trains a model utilizing its dataset and subsequently shares the model parameters with other participants. Crucially, the actual data does not leave the local premises, thereby maintaining the confidentiality of sensitive patient information. This approach facilitates collaborative learning without compromising data privacy.
% the distributed computation and develpmoent of deep learning models ensure the accumulation of knowledge and reaching the full capacity of deep learning models.

FL can differ from centralized data sharing in a number of ways. While both approaches aim to optimize their learning objective, FL algorithms have to account for the fact that communication with clients takes place over unreliable networks with very limited upload speeds. So unlike 
the centralized setting in which computation is generally a bottleneck, in FL communication might be the bottleneck. 

In this study, a framework was developed to facilitate collaboration among hospitals by employing multiple data sources for the detection of COVID-19 infections via FL. The framework's decentralized data distribution ensures privacy, as data remains stored locally\cite{darzidehkalanifederatedII}.




% There are several ways of doing federated learning since federated learning is performed in a higher level compared to conventional deep learning. There are more parameters of complexity since the model should take into consideration the general and local methods of optimization, data accumulation for all the sites, and local training models. 
% There are several parameters playing a role in the final performance of a federated learning setting.

 





% It can be spread by human-human interaction. There are several waus tp diagnose COVID-19. Transcription-polymerase chain reaction (RT-PCR) test are one of the most common ways of diagnosing COVID-19. However, there are other ways of diagnosing COVID-19 \cite{rubin2020role}
% However, due to large number of people having the symptoms and being tested for COVID-19, it might not always be feasible for radiologists to examine large number of patients during the outbreak. Scans of CT and Xray of infected patients have sorts of deformation, misalignment between certain parts, or have pixel intensities different than normal people in infected parts. Observing this misalignment in not an easy task and requires effort and focus from radiologist to correctly distringuish betweeen healthy and infected people. 
% Instead, 
% Various architectures of CNN have been tested  and shown a great performance on the imaging datasets. Models such as AlexNet \cite{krizhevsky2012imagenet}
%  ResNet \cite{he2016deep}  and MobileNet \cite{howard2017mobilenets} could classify thousands of objects belonging to hundreds of classes with the same accuracy as human.
 



% It has also been used to detect COVID-19 from imaging data. Most of the usecases were either for Xray like the researches in 
% While CT scans can help screen COVID-19 potential infected people, CT scans of other diseases could also have overlapping features with COVID ingected patients. Hence, it is tricky to distinguish COVID infected patients from other types of lung inflammation and requires radiologists level experrtise. The most literature focuses on the ways to better distirnguish COVID infected patients from other diseases, as can be seen in [8], [9], [10], [11], [12]. 

\section{Background and Related works}

 % $\Delta W_n^t$
 
Federated learning has demonstrated efficacy in an array of imaging modalities, including Magnetic Resonance Imaging (MRI) \cite{sheller2020federated}\cite{silva2019federated}, X-ray \cite{balachandar2020accounting}, retinal imaging \cite{balachandar2020accounting}, as well as in applications such as brain tumor segmentation \cite{bakas2017advancing}\cite{lee2018privacy}, diagnosis \cite{pan2019improving}, and treatment selection \cite{lee2018privacy}. In particular, Federated Learning (FL) has proven to be a valuable tool for supporting physicians in their decision-making process regarding the treatment of COVID-19 patients. A landmark study that involved 20 institutions across five continents found that FL played a significant role in shaping patient treatment plans\cite{flores2021federated}. The study employed chest radiography images in conjunction with clinical information to determine the appropriate level of care and oxygen requirements for patients afflicted with COVID-19. It was observed that FL improved the performance of the predictive model, especially for institutions with smaller datasets, compared to using only local data for model training. Additionally, it was found that healthcare facilities with smaller datasets often had underrepresented categories due to a low number of patients in certain classes. The implementation of FL led to a notable improvement in predictions for these underrepresented patient categories.

\begin{figure*}[t!]
\centering
\includegraphics[width=1\textwidth]{FLmodels.png}
\caption{Illustration of FL models and algorithms: (a) Federated averaging, where clients train on a local batch of data. (b) FedSGD, in which a subset of clients is selected, and each performs a single step of SGD before sending model updates to the server. (c) Cyclic Weight Transfer (CWT), where clients train locally and pass the model to the next client, repeating the cycle. (d) Single Weight Transfer (SWT), where the model passes through each client only once. (e) Stochastic Weight Transfer (STWT), in which the model is sequentially passed through clients, with participating clients in each round being sampled randomly.}
\label{fig:flalgorithms}
\end{figure*}


Recent research has focused on the classification of scan images to distinguish between COVID-19 patients and healthy individuals, as well as identifying lesion areas. The primary application of AI in managing COVID-19 patients has been the interpretation of radiology images, especially chest CT scans.  The detection of lung alterations through these scans plays an important role in optimizing patient management and guiding treatment decisions\cite{yan2020interpretable}\cite{hu2020challenges}\cite{burian2020intensive}.  Several studies have also explored 3D Convolutional neural networks \cite{wang2020weakly} and COVID-19 detection with a limited number of training samples.

While the majority of these studies report favorable accuracy, they often presume a centralized environment wherein a single data center has access to all data. However, a few studies have successfully applied distributed learning for COVID-19 detection, employing global aggregation models such as model averaging in federated learning settings\cite{ho2022fedsgdcovid}\cite{zhang2021dynamic}, or within a blockchain framework \cite{kumar2021blockchain}. These studies have pointed to certain limitations of existing algorithms, such as high communication overhead \cite{remedios2020federated}, as well as convergence issues or catastrophic forgetting when the number of participating hospitals increases \cite{sheller2020federated} \cite{chang2018distributed}.


To our knowledge, no study has been performed that compared multiple FL algorithms under standard conditions to evaluate their applicability. Therefore, comparing multiple FL algorithms under standard conditions could be informative in evaluating their applicability in practice.

To evaluate the existing methods from multiple perspectives, we have implemented the most popular models and compared them in terms of performance, communication overhead, and computation burden.


\section{Algorithms}
% \label{sec:algorithms}
\textbf{Centralized Data Sharing} In Centralized Data Sharing (CDS), data is stored in a central location and is accessible to all clients. This stands in contrast to federated and decentralized data sharing methods, where data is stored across multiple locations and accessed by either a single user or a limited number of users. CDS serves as a baseline for comparing other algorithms.
\textbf{Federated Averaging} Federated Averaging involves an iterative learning procedure comprising local and global steps. In this process, each data owner trains a model received from a global server on its local dataset through local iterations \cite{mcmahan2017communication}. Subsequently, the global server aggregates the updated local models to update the global model. This global model is then distributed to clients for the subsequent round. The optimization problem for Federated Averaging can be expressed as:
\begin{equation}
w^{t+1} = \sum\limits_{i=1}^{N}{p_{i} w_{i}^{t}} , w_{i}^{t}=\arg\min\limits_{w_{i}}{\left(\mathcal{L}(\mathcal{D}_{i};w^{t})\right)}
\end{equation}
where $N$ is the number of data owners, $\mathcal{L}(\mathcal{D}_{i};w^{t})$ is a loss function indicating global model parameters  $w^{t}$ of local datasets, and $p_{i}$ is the probability of selecting client $i$. 
Local optimization can be formulated as $w_{i}^{t+1} \leftarrow w^{t}-\eta\cdot \nabla \mathcal{L}(w^{t};\mathcal{D}_{i})$, where 
 $\eta$ is the learning rate. The global model can be updated based on the local models $w_{i}$ and is shared for aggregation: 
\begin{equation}
w^{t+1} = \sum\limits_{i=1}^{N}{p_{i} w_{i}^{t+1}}
\end{equation}

\textbf{Federated Stochastic Gradient Descent} Federated Stochastic Gradient Descent, or FedSGD\cite{chai2020fedeval}, is a variant of Federated Averaging (FedAvg) that employs a large-batch synchronous approach for multi-client learning. In FedSGD, a subset of clients is chosen from the total pool, with this subset being defined by $C$. This selection is made at each global round. Subsequently, the global server dispatches the latest global model to the chosen clients. Each of these clients then undertakes local training on its dataset for a predetermined number of epochs. The global model is updated in light of the local models obtained from each client and is made available for aggregation, a process that is analogous to that in FedAvg. Nevertheless, FedSGD distinguishes itself by computing the gradient over the chosen batch of clients, hence $C < 1$. When $C = 1$, the training becomes non-stochastic, or full batch, as it involves all the clients. This makes it possible to train with large batches since the gradient is computed over the selected subset of clients.

The optimization problem for FedSGD can be expressed as follows:
\begin{equation}
w^{t+1} = w^{t} - \eta \cdot \sum\limits_{i=1}^{C}p_i \cdot \nabla \mathcal{L} (w^t;\mathcal{D}_i)
\end{equation}
where $\eta$ is the learning rate, $p_i$ is the probability of selecting client $i$ and $\mathcal{L}$ is the loss function.
The key difference between FedAvg and FedSGD lies in the use of large-batch synchronous approach in FedSGD. This approach has been shown to outperform the naive asynchronous SGD training due to the increased accuracy and efficiency, as compared to the local training approach used in FedAvg \cite{chai2020fedeval}\cite{charles2021large}. Additionally, FedSGD has been shown to be more robust to non-IID data distributions, compared to FedAvg \cite{chai2020fedeval}. 
\quickthings{This part is unclear to me.}
% \absolutelynecessary{}
\maybelater{paraphrised}

\textbf{Cyclic weight transfer} Federated learning techniques have been widely used in medical image processing tasks using a method known as cyclic weight transfer (CWT)\cite{balachandar2020accounting}. This method involves training models on individual clients for a number of iterations and then cyclically sharing the updated weights with the following client. However, the existing CWT algorithm faces a notable challenge, as it lacks the ability to effectively manage inter-client variability in training data or labels.  To ensure the practical application of CWT, it is crucial to develop a version that can handle the common variations observed in a majority of real-world medical imaging datasets\cite{darzidehkalanifederatedII}.

\textbf{Single weight transfer} Single weight transfer (SWT) is anothe FL method widely used in the medical imaging domain. In Single weight transfer, models are trained in each client with its local data, and then the updated model is transferred to the next client. The difference between this method and CWT is that here the model passes each client only once\cite{darzidehkalanifederatedI}. 

\textbf{Stochastic weight transfer}
In stochastic weight transfer (STWT), we select a subsample of clients and train them in a cyclic manner. Similar to FedSGD, a ratio defines the number of selected clients to the total number of clients in each federated round. 

Figure \ref{fig:flalgorithms} provides a visual representation of the introduced algorithms.



% سوال / مسیله چیه
% Importance: Why your research matters in the context of an industry or the world
% اهمیت موضوع
% 	کجاها کاربرد داره
% مشکلات موضوعات قبلی و اینکه مسایل دگ کار نمیکنن
% برای اینکه این مسایل کار کنن نیازی به چه چیزی بود. ک اونا کم داشتن
% چرا این روش تو داره کار میکنه
% چه سوالی رو جواب میده
% توضیح اینکه چجوری کار میکنه
% تعریف و تمجید از کار کردنش

% \maybelater{chanta akse tamiz ba coreldraw bekesh}

% Their
% We are comparing various implementations of FL and comparing their performance level. The performance is done by metrics.
% 2. Although FL has been introduced to tackle the problem of privacy, its distributed nature requires too much commuincation between clients. Knowsing can determine its deployability and scailibitlty. Hence the question is how to these model compare it terms of communication between clients.
% 3. Computation time. Each FL model is investigated based on the required time for computation.
% 4. Effect of rounds: Number of Federated rounds 


% \quickthings{Maybe you can add a benchmarking section similar to "A Performance Evaluation of FL Algorithms"}
% \maybelater{Some ideas are herer https://arxiv.org/pdf/1709.05929.pdf
% }





\section{Experiment}




\subsection{Datasets}

\textbf{Brain cancer classification:}Dataset for detecting brain cancer is downloaded from Kaggle. It contains brain MRI images categorized into four classes: Three types of brain tumor,  and one healthy no Tumor class. 
% For tumor type classification task, 2870 images for training and 394 images for the test. 
For Meningioma and Glioma detection, 1437 and 1426 samples were used, respectively,  and they were split as train/ test data. Transforms rotation, Flipping, and normalizing.  Images were resized to 100$\times$100. 
\\
\textbf{Histopathologic cancer detection:} 
The samples are metastasic tissue images of lymph node cancer. Images are in size ,96 $\times$ 96, and the task is to classify the tissue samples. Each positive sample has a metastasic region, which is located in $32\times32$ neighborhood of each image sample. 
% Tumor tissue in the outer region of the patch does not influence the label. 
They were categorized into cancer and non-cancer classes. 2150 samples were chosen and randomly assigned to clients, and each split 62\% for train and 38\% for the test. We used horizontal Flipping, and images were normalized.

% \\\hl{A ro tu line 21 algorithm taarif kon va dorost kon notionesh ro}
\subsection{Network architectures}

\textbf{Deep learning models:}
The deep learning model used is a Convolutional Neural Network (CNN) with six layers of convolution stacked to 5 fully connected layers. The activation function used is ReLU. And dropout parameter (0.25). The last layer depends on the number of classes (4 or 2). Models are all trained and converged before implementing adversarial attacks.
They are trained with Cross-Entropy loss and with an SGD optimizer. Four models are trained separately. Two binary classification models are trained to detect Meningioma and Glioma from the healthy No Tumor class. One model is trained to detect cancer in HistoPathology images.

\textbf{Federated setting:}
Three clients are defined for the FL setup. The data is assigned randomly, and the clients have non-IID data distribution.
The FedAVG method is used to aggregate the models. The aggregation is weighed based on the length of the training dataset. Each client is trained for 20 epochs at the communication round. The total FL rounds are 50.

\subsection{Attack setting}

Each client is assigned 100 test samples to perform the adversarial attack. Each client is selected as the adversary, and the results are averaged at the end. 
% The attacks are conducted with different $\epsilon$ and gradient steps $\alpha$ to evaluate the effect on the final result. 
We use three metrics,
Clean accuracy is defined as the performance of models on uncorrupted test images, \textit{Attack Success Rate (ASR)}; how much an adversary can change the predicted labels produced by each model, \textit{Average Error Transferability Rate (AETR)} to evaluate transferability. An adversary might propose an attack by only selecting the samples that successfully fooled its own model. AETR measures how successfully these samples changed the target models' correct predictions. So unlike ASR, this metric doesn't count the examples that were initially misclassified by the victim model.\cite{papernot2017practical},

\begin{figure*}[h!]
     \centering
     \begin{subfigure}
         \centering
         \includegraphics[width=0.32\textwidth]{Meningioma_ASR_EPS.eps}
          \includegraphics[width=0.32\textwidth]{Glioma_ASR_EPS.eps}
           \includegraphics[width=0.32\textwidth]{Pathology_ASR_EPS.eps}
         \caption{Effect of $\epsilon$ on Average ASR on benign clients }
         \label{fig:AASR benign-EPS}
     \end{subfigure}
     \hfill
     \begin{subfigure}
          \centering
         \includegraphics[width=0.32\textwidth]{Meningioma_others_EPS.eps}
          \includegraphics[width=0.32\textwidth]{Glioma_others_EPS.eps}
           \includegraphics[width=0.32\textwidth]{Pathology_others_EPS.eps}
         \caption{Effect of $\epsilon$ on ASR on the advesarial client }
         \label{fig:AASR malignant-EPS}
     \end{subfigure}
     \hfill
        \caption{Effect of Error perturbation degree $\epsilon$ on attack transferability. FGSM and PGD attack with and without CRN initialization where performed. CRN initalizations used 10 models and another one loading 5.  ASR is calculated on benign and adversary clients. The results of with varying perturbation degrees. The higher ASR on benign clients shows higher transferability}
        \label{fig: epsilon graphs}
\end{figure*}

  Here, we expand the previous findings to the FL settings and discuss whether change in $\epsilon$ can lead to transferability. 
  
We perform three analysis to find important factors in attack success:
\begin{itemize}
    \item We investigate dependency of attacker on $\epsilon$. By visual inspection and ASR evaluation, and as a comparison to previously found optimal known values in centralized  setting.
    \item  We compare  efficiency of models, by discussing their attack preparation time and their ASR.
    \item We also see how attack step $\alpha$ can determine ASR.
\end{itemize}



\section{Results}







% Original images, adversarial images, and corresponding adversarial noise created with FGSM ( ϵ=0.02) in different black-box settings


\subsection{Efficiency analysis}


\begin{figure*}[h!]
     \centering
     \begin{subfigure}
         \centering
         \includegraphics[width=0.32\textwidth]{Meningioma_ASR_EPS_steps.eps}
          \includegraphics[width=0.32\textwidth]{Glioma_ASR_EPS_steps.eps}
           \includegraphics[width=0.32\textwidth]{Pathology_ASR_EPS_steps.eps}
         \caption{Average ASR on benign clients }
         \label{fig:y equals x}
     \end{subfigure}
     \hfill
     \begin{subfigure}
          \centering
         \includegraphics[width=0.32\textwidth]{Meningioma_others_EPS_steps.eps}
          \includegraphics[width=0.32\textwidth]{Glioma_others_EPS_steps.eps}
           \includegraphics[width=0.32\textwidth]{Pathology_others_EPS_steps.eps}
         \caption{Average ASR on advesarial client }
         \label{fig:y equals x}
     \end{subfigure}
     \hfill
        \caption{Effect of Error perturbation step $\alpha$ on attack transferability. PGD attack with CRN initalizations used 10 models and another one loading 5. ASR is calculated on benign and adversary clients. The results of with varying perturbation steps. The higher ASR on benign clients shows higher transferability}
        \label{fig:three graphs}
\end{figure*}

The models are trained in an FL enviroment. Average test results, or clean accuracy shows the average performance of clients on unperturbed test data. This could be a baseline for comparison of attacks.
Here we compare the computational complexity of iterative algorithms. The comparison is based on the time required to train one batch of adversarial examples, with  $\epsilon=0.05$ and $\alpha=0.001$. The average Attack success rate is the average of ASR for all test samples of all clients and indicates the transferability in the FL network.
As expected \ref{table_time}, the higher iteration rounds take more computation and lead to higher AASR. Generally, computational complexity has a linear dependence on the number of iterations. However, the increase in transferability can vary. If PGD performs poorly initially, having higher iterations doesn't help much. 

Also, CRN does not require additional computation. Similar to single-step methods, computational load is small but increases on high-resolution images. As an example, the pathology dataset requires more calculation than MRI data. 
However, the overall time is far less than iterative models, and also it could enhance the transferability much better. The effect is consistent among different datasets. 
% However,  This means that Pathology images have better AASR than MRI images in baseline and CRN enabled attacks.
% Also,  

% PGD can sometimes perform poorly, even in higher iterations. Enabling the initialization might  


\begin{table}[h!]
\centering
\setlength{\tabcolsep}{6pt}
\renewcommand\arraystretch{1.22}
\caption{ \small Comparison of iterative models and their copmutational efficiency, on performing computations on one batch of data.  ACC shows average client performance for unperturbed test data. AASR is average ASR on all clients.}
\begin{tabular}{| *{5}{c|} }
\hline
Dataset  & Attack type & ACC & AASR & time (sec)
\\   \hline  
\multirow{3}{5em}{Meningioma}     &\textbf{PGD-1+CRN10}&\multirow{3}{3em}{84.12\%}&\textbf{63.92\%} & \textbf{0.217}  \\
 &PGD-20&&27.52\% & 3.423 \\
&PGD-40&&32.99\% & 6.794  \\ \hline
\multirow{3}{4em}{Pathology}     &\textbf{PGD-1+CRN10}&\multirow{3}{3em}{77.01\%}&\textbf{90.92\%} & \textbf{0.354} \\
&PGD-20&&60.98\% & 5.464\\
&PGD-40&&82.27\% & 10.843 \\ \hline
\multirow{3}{3em}{Glioma}     &\textbf{PGD-1+CRN10}&\multirow{3}{3em}{61.84\%}&\textbf{70.55}\% & \textbf{0.216}  \\
&PGD-20&&51.83\% &3.420  \\
&PGD-40&&63.77\% & 6.793  \\ \hline


\end{tabular}
\label{table_time} 
\end{table}




\subsection{Effect perturbation degree}

We performed attack scenarios to evaluate the effect of perturbation degree $\epsilon$  on attack success on the model. We compared the baseline attack to CRN- enabled attacks. ASR is calculated on benign and adversary clients and is shown in Fig \ref{fig: epsilon graphs}.

More successful attacks on adversaries led to higher transferability to benign clients. Low values of $\epsilon$ can decrease transferability to a large extent. $\epsilon$ = 0.01 has poor performance on all the clients. Generally, Higher $\epsilon$ values lead to higher ASR in all scenarios. Although for Meningioma and Glioma images, increasing $\epsilon$ from 0.04 didn't have much positive effect in PGD and FGSM methods on the adversarial client. 
In pathology images, Iterative models can reach perfect accuracy on the adversarial client and have very high transferability.

CRN  has a positive effect on transferability in all scenarios. Also, CRN-enabled models are more dependent on $\epsilon$; their extent can vary depending on the dataset, attack method, and value of $\epsilon$.
Tuning $\epsilon$  might be tricky since performance and imperceptibility should be considered together.High $\epsilon$ values mean more distortion in the image. We performed an FGSM attack with varying $\epsilon$ values to evaluate the perturbation effect visually. Fig \ref{fig:perturbedMRI} shows the results of perturbed images with different perturbation degrees and their comparison with the unperturbed image. Note that perturbation values below  $\epsilon=0.05$ cause limited distortion.


Iterative models can reach high ASR very quickly with increasing $\epsilon$. 

\begin{figure*}[h!]
 \centering
 \includegraphics[width=0.8\textwidth]{Adversarial_attacks-5.eps}
 \caption{\small From (a) to (f), normal tumor image, and petrurbed images with FGSM attack, with perturbation parameters: $\epsilon$ = 0.01, 0.03, 0.05,  0.10, respectively}
 \label{fig:perturbedMRI}
\end{figure*}











% \subsection{AATR}
% \begin{table}[h]
% \centering
% \setlength{\tabcolsep}{8pt}
% \renewcommand\arraystretch{1.4}
% \caption{Result of our experiments for PGD attack on Pathology data $\epsilon$ values. The iterative models are trained for 40 rounds with step size of 0.01. \hl{\faQuestion maybe you can change the word for without CRN}}
% % \begin{tabular}{cccccccccccc}
% \begin{tabular}{| *{4}{c|} }\hline
% % Dataset  & Attack type & \multicolumn{2}{c|}{ $\epsilon=0.01$} & \multicolumn{2}{c|}{ $\epsilon=0.03$}           &    \multicolumn{2}{c|}{ $\epsilon=0.05$}    & \multicolumn{2}{c|}{ $\epsilon=0.07$}                                                              \\ \hline
%  Dataset&\multicolumn{3}{c|}{ FGSM}    
% %  &              \multicolumn{3}{c|}{ PGD}                                                                           \\  

%  \cline{2-5}
% &Baseline&CRN5 &CRN10\\\hline

% Meningioma&84.80\%&\textbf{86.31}\%&{84.02\%}
% % \%&83.73\%  &\textbf{84.25}\%  &84.16\%
% \\\hline
% Glioma&82.32\% &87.06\%&\textbf{89.98\%}
% % &74.89\%&\textbf{83.64}\% &82.32\%
% \\\hline

% Pathology&\textbf{90.67\%}&{84.81\%}\%&{85.98\%}
% % \%&79.86\% &86.79\%&\textbf{86.93\%}
% \\\hline



% \end{tabular}
% \label{table_core} 
% \end{table}










\subsection{Effect perturbation step}


Another critical parameter is the step of change in iterative methods. {$\alpha$} is shown to be highly deterministic in specific settings, but is not yet investigated in MIA domain. To evaluate how $\alpha$ can affect transferability, we examined the clients with $\eps=0.03$. Under various steps. The general observation is that there is an optimal middle range the same as $\epsilon$. The step size of 0.007 led to the highest transferability in MRI images. Larger steps led to a sharp decrease in transferability. For pathology images, midspans, about 0.007 led to the best ASR on the adversary, and higher values decreased both ASR on adversary and transferability. However, smaller steps also had high ASR values. 







\subsection{Error tranferability}




Another measure is the average error transfer rate (AETR). Although PGD outperforms FGSM in ASR, in AETR, the results are close. PGD is trained for 40 iterations, and the models are compared with and without CRN.
Both FGSM and PGD having high AETR suggest that although FGSM might not be able to produce high-quality samples in general, the good examples which can fool the adversary are also highly transferable to benign clients. And can be better than PGD samples. CRN generally increases the AETR. 




\begin{table*}[h!]
\centering
\setlength{\tabcolsep}{8pt}
\renewcommand\arraystretch{1.4}
\caption{Result of Average Error Transfer Rate (AETR) for FGSM and PGD methods, with and without CRN initalization }
% \begin{tabular}{cccccccccccc}
\begin{tabular}{| *{7}{c|} }\hline
% Dataset  & Attack type & \multicolumn{2}{c|}{ $\epsilon=0.01$} & \multicolumn{2}{c|}{ $\epsilon=0.03$}           &    \multicolumn{2}{c|}{ $\epsilon=0.05$}    & \multicolumn{2}{c|}{ $\epsilon=0.07$}                                                              \\ \hline
 Dataset&\multicolumn{3}{c|}{ FGSM}     &              \multicolumn{3}{c|}{ PGD}                                                                           \\  

 \cline{2-7}
&Baseline&CRN5 &CRN10&Baeline&CRN5&CRN10\\\hline

Meningioma&84.80\%&\textbf{86.31}\%&{84.02}\%&83.73\%  &\textbf{84.25}\%  &84.16\%\\\hline
Glioma&82.32\% &87.06&\textbf{89.98}&74.89\%&\textbf{83.64}\% &82.32\%\\\hline

Pathology&\textbf{90.67\%}&{84.81}\%&{85.98}\%&79.86\% &86.79\%&\textbf{86.93\%}\\\hline



\end{tabular}
\label{table_core} 
\end{table*}


% \newpage


\section{Discussion}
\label{sec:discussion}

% Local training had near random accuracy, but in all the experiments, the performance results were improved. 
% Considering low number of samples for individual clients, a training process limited to few clients results in near random classification accuracy. 
% Local models can have low performance due to limited datasets.
% Also,, ImageNet pre-training had a speed-up effect on our model performance.

% \textbf{Fundamental groups of FL}

% \textbf{how many groups}
% \textbf{how can we group them, name different ways}
% \textbf{effect number of data}


% Paragraph 1: This paragraph provides a “big picture” perspective for readers to remind them of the importance of your study.

% \textbf{overall results}
% 
Our results show that FL has comparable performance to centralized data sharing, with the advantage of keeping data private. With large volumes of data and after high number of rounds, centralized data sharing and cyclic weight transfer have the highest accuracy. 


Sequential models are susceptible to catastrophic forgetting, where a global model performs well on the latest client it has seen while having poor performance in other clients.\quickthings{cryptic sentence}\maybelater{I have paraphrised the sentence}
Conversely, in algorithms like FedAvg and FedSGD, the models are averaged asynchronously after all the clients have finished their training. So the trajectory is smoother and overall improving with more communication rounds. As shown in Figure \ref{seq-vs-noseq}, local test results can have a high variance when passing through clients sequentially, indicating the catastrophic forgetting effect.

Models like FedAvg, and FedSGD, in which all the clients have identical copies of one global model, are slower and more challenging to converge compared to sequential models like CWT and STWT. Also, FedAvG and FedSGD require more training resources due to active server participation, resulting in more computation and network consumption. Stochastic client selection is an efficient way of training. Stochastic models save significant time and resources while having similar performance to full client participation. Overall, CWT and STWT have best results in terms of model accuracy and computation times. These findings could be practical in further federated deployments in medical institutions.
% but when there are communication and computation limitations, other methods might be more practical.




% As can be seen in the results, CWT and STWT have better local results than other methods.

% So each local model is different than other models, and even in fewer rounds they retrain the model on their local data and achieve high accuracies. 

%  Unsurprisingly, CDS has higher accuracy than non-sequential methods.
 
% so in limited number of rounds, latter algorithms are preferred.



%  FedAVG needs to perform both local training and global aggregation, which requires is to have more computation, than sequential models which eliminate global aggregation and broadcasting step.   
%  Models with global server, suh as FedAVG and FedSGD, required more data transfer, because of the two way communication of each client with server.
Sequential models like CWT and STWT perform better than non-sequential models on fewer training rounds. For example,  after three rounds of training, STWT and CWT both reach 96\% accuracy, while FedAvG reaches 66\%, and FedSGD performs equally to a random classifier. As the training proceeds, FedAVG and FedSGD gradually improve with more global rounds.The concept of sequential models is similar to fine-tuning \cite{chen2020online} in centralized deep learning, so in cases where a hospital temporarily joins an FL network, or there is an urgency in training, sequential models are a better option.

More training rounds do not always lead to a better global model. Although average performance on all clients improves, more global rounds lead to worse performance for some clients. The global model can overfit some clients, leading to lower performance on others\cite{mohri2019agnostic}. Some studies suggested early stopping and fine-tuning to local dataset after global training is finished \cite{yu2020salvaging}. 
In all the algorithms, more clients resulted in slower convergence. This effect is stronger in the FedAvg algorithm. In FedAVG, the Global model must compromise between potentially disparate local minima\cite{li2019convergence}. Methods such as adaptive or stochastic selection of clients and momentum-based models help faster convergence \cite{liu2020accelerating}.
Our results suggest that stochastic client participation is close to full client participation. The average results of four trials with varying rounds, shown in Table \ref{number_of_rounds} indicate that stochastic client participation in FedSGD results in 5.23\% performance loss and 40\% less bandwidth consumption compared to FedAvg. In STWT, it results in only 1.25\% less accuracy but saves 40\% of communication and 11.3\% of computation.
% Also, SWT improves local clients' performance up to 80\% accuracy with 8 clients, and it requires extremely low bandwidth requirements.
These results are in accordance with prior studies, showing that, in theory, stochastic and full client participation have similar global minima\cite{cho2020client}. Stochastic client selection can be advantageous when there are limited resources, or in larger networks with occasionally unavailable clients.


% Models like FedSgd and STWT, select the participating clients in a stochastic manner. 
% The reason is because each client has less data, and there will be more number distant local client minima leading to unstable global model.


% Such can be a reflection of real world setting, where not all the clients are active in each round. So a fraction are selected.


% Some models have high max test accuracy while average is low indicating bias and catastrophic forgetting. Infrastructure limitations should be considered, when a central server has limited computing power it would be unable to perform FedAVG, instead no-server algorithms require to pass the model to the next client.

We did not assume any shift in clients' data. A more comprehensive analysis should consider the effect of the domain and distribution shifts on the performance of the algorithms. Also, inter-client data variability and the effect of heterogenous clients could be a future line of research.
% For deployments in longer timeframes, there might be changes in clients data distirbution , and previous results might be inapplicable Effect of domain shifts, and inter-client data variability could be a future line of research. 






\section{Conclusion}
\label{sec:conclusion}
 
FL enables extensive collaborations of hospitals to address medical imaging problems while keeping data private. Real-world implementation requires consideration of efficiency and hardware requirements in addition to model performance, especially in the healthcare field, which generally has limited infrastructure. We implemented five FL algorithms for COVID-19 detection and analyzed their efficiency and accuracy.
Our results suggest that FL algorithms have comparable performance to centralized data sharing, with the advantage of keeping data private. They also show that the sequential methods are a better option in most of the scenarios. This study can be helpful in the deployment of FL systems in COVID-19 detection and medical image analysis in general.

% demonstrate better performance for detecting COVID-19 patients and might be practical in deploying FL algorithms for covid-19 detection and medical image analysis in general.





\printbibliography






% \section*{Acknowledgement}



\ifCLASSOPTIONcaptionsoff
  \newpage
\fi



\bibliographystyle{IEEEtran}
\bibliography{IEEEabrv}



\vfill
% \section{To do list}
% \subsection{Quick things}

% \subsection{Absolutely necessary}

% %   \item add pathology image to figure perturbation of MRI
% \subsubsection{Paper}
% \begin{itemize}
%     % \item paraphrise + shorten FL and DP sections
%     % \item baad yebar ba dide Chaoning bekhun
% \end{itemize}
% \subsubsection{Coding}
% \begin{itemize}

% \end{itemize}
% \subsection{Maybe later}
% \subsubsection{coding}
% \maybelater{CWT but with random ordering}
% \maybelater{One shot FL?}
% \maybelater{Vazne kamtari dadan be latest client in CWT ?}
% \maybelater{Testing with other deep networks?}
% \maybelater{At least you can increase depth of network}
% \maybelater{Impact of data heterogeneity? just split the data heterogeneously and see how do they work?}

% \maybelater{Effect of Differential privacy?}
% \maybelater{Effect of pretrained=False?}
% \maybelater{HasSpike model?}
% \maybelater{Maybe small variations to generate new models}
% \maybelater{Maybe doing gradient masked FedAvG?}
% \maybelater{Maybe evaluating different number of local epochs like https://github.com/siddarth-c/FedGMA}
% \maybelater{Adding all these ideas to Obsidian/presentations}
% \maybelater{Going for another paper if an idea is worth pursuing}
% \maybelater{Another paper: CWT/SWT with DP}
% \maybelater{Local GD like Ahmed khaled}

% \subsubsection{Paper}
% \begin{itemize}
%     \maybelater{Adding mathematical expression of each FL model (good to improve math level of paper)}
%     \maybelater{maybe adding detailed graph for each participant as an Appendix?}
%     \maybelater{Adding algorithmic view of each FL alg?}
  
% \end{itemize}





% that's all folks
\end{document}


